\documentclass[12pt]{article}
\usepackage{amsmath, amssymb}
\usepackage[margin=1in]{geometry}
\setlength{\parskip}{6pt}
\title{QUBO Formulation: Minimum Graph Bisection Problem}
\date{}
\begin{document}
\maketitle

\subsection*{Minimum Graph Bisection Problem}

\paragraph{Problem Statement.}
Given an undirected weighted graph $G = (V, E)$ with an even number of vertices $N = |V|$ and non-negative edge weights $w_{ij}$ for each $(i,j) \in E$, partition the vertex set $V$ into two disjoint subsets $V_A$ and $V_B$ such that $|V_A| = |V_B| = N/2$. The objective is to minimize the total weight of edges crossing between the two partitions, i.e., minimize $\sum_{(i,j) \in E} w_{ij} \mathbb{I}(i \in V_A, j \in V_B \text{ or } i \in V_B, j \in V_A)$.

\paragraph{Decision Variables.}
For each vertex $i \in V$, introduce a binary variable $x_i \in \{0, 1\}$. The assignment $x_i = 0$ indicates that vertex $i$ belongs to partition $V_A$, while $x_i = 1$ indicates membership in partition $V_B$. The set of all variables is $\{x_i\}_{i \in V}$.

\paragraph{QUBO Derivation.}
\textbf{Step 1 --- Objective Formulation.}
The cut weight can be expressed using the binary variables as $\sum_{(i,j) \in E} w_{ij} (x_i - x_j)^2$, since $(x_i - x_j)^2 = 1$ when $x_i \neq x_j$ and $0$ otherwise. For binary variables, the identity $(x_i - x_j)^2 = x_i + x_j - 2x_i x_j$ holds. Thus, the minimization objective is:
\begin{align}
    \min \sum_{(i,j) \in E} w_{ij} (x_i + x_j - 2x_i x_j).
\end{align}

\textbf{Step 2 --- Constraint Formulation.}
The requirement that both partitions have equal size is enforced by the equality constraint:
\begin{align}
    \sum_{i \in V} x_i = \frac{N}{2}.
\end{align}

\textbf{Step 3 --- Penalty Term Expansion.}
We convert the equality constraint into a quadratic penalty term with parameter $\lambda > 0$:
\begin{align}
    P(x) = \lambda \left( \sum_{i \in V} x_i - \frac{N}{2} \right)^2.
\end{align}
Expanding the square yields:
\begin{align}
    P(x) = \lambda \left( \sum_{i \in V} x_i^2 - N \sum_{i \in V} x_i + \frac{N^2}{4} + \sum_{i \neq j} x_i x_j \right).
\end{align}
Using the idempotent property of binary variables, $x_i^2 = x_i$, and noting that $\sum_{i \neq j} x_i x_j = 2 \sum_{i < j} x_i x_j$, we obtain:
\begin{align}
    P(x) = \lambda \left( (1 - N) \sum_{i \in V} x_i + 2 \sum_{i < j} x_i x_j + \frac{N^2}{4} \right).
\end{align}

\textbf{Step 4 --- Combined Hamiltonian.}
Adding the objective and penalty terms gives the full QUBO Hamiltonian $H(x)$:
\begin{align}
    H(x) = \sum_{(i,j) \in E} w_{ij} (x_i + x_j - 2x_i x_j) + \lambda \left( (1 - N) \sum_{i \in V} x_i + 2 \sum_{i < j} x_i x_j + \frac{N^2}{4} \right).
\end{align}
Grouping linear and quadratic terms, and defining the weighted degree $\deg_w(i) = \sum_{j: (i,j) \in E} w_{ij}$, we rewrite $H(x)$ as:
\begin{align}
    H(x) = \sum_{i \in V} \left( \deg_w(i) + \lambda(1 - N) \right) x_i + \sum_{(i,j) \in E} (-2w_{ij}) x_i x_j + \sum_{i < j} (2\lambda) x_i x_j + \frac{\lambda N^2}{4}.
\end{align}

\textbf{Step 5 --- Q Matrix and Offset.}
The Hamiltonian takes the standard QUBO form $H(x) = x^T Q x + c$. The matrix entries and scalar offset are:
\begin{align}
    Q_{i,i} &= \deg_w(i) + \lambda(1 - N), \quad \forall i \in V, \\
    Q_{i,j} &= \begin{cases} 
        -2w_{ij} + 2\lambda & \text{if } (i,j) \in E, \\
        2\lambda & \text{if } i \neq j \text{ and } (i,j) \notin E, \\
        0 & \text{otherwise},
    \end{cases} \\
    c &= \frac{\lambda N^2}{4}.
\end{align}

\paragraph{Penalty Weight Justification.}
To guarantee that the global minimum of $H(x)$ over $\{0,1\}^N$ corresponds to a feasible bisection, the penalty parameter $\lambda$ must dominate any potential reduction in the cut weight achievable by violating the balance constraint. Flipping a single variable changes the constraint violation by exactly $1$, increasing the penalty by $\lambda$. The maximum possible decrease in the objective function from a single bit flip is bounded by the maximum weighted degree $\Delta_w = \max_{i \in V} \deg_w(i)$. Setting $\lambda > \frac{\Delta_w}{N+1}$ ensures that the energy cost of a unit constraint violation strictly exceeds any feasible objective improvement, thereby forcing the optimizer toward feasible configurations.

\paragraph{Correctness Argument.}
Minimizing $H(x)$ over binary vectors recovers the optimal bisection because (i) all feasible assignments satisfying $\sum_{i \in V} x_i = N/2$ incur zero penalty, reducing $H(x)$ exactly to the original cut weight objective; and (ii) any infeasible assignment incurs a penalty strictly greater than the maximum possible decrease in the cut weight from any single variable flip. Consequently, the global minimum of $H(x)$ must occur at a feasible point, and among feasible points, $H(x)$ is identical to the cut weight. Thus, the minimizer of $H(x)$ corresponds precisely to the minimum-weight balanced partition.

\paragraph{Illustrative Example.}
Consider a path graph with $N=4$ vertices $V=\{0,1,2,3\}$ and edges $E=\{(0,1), (1,2), (2,3)\}$ with weights $w_{01}=1$, $w_{12}=2$, $w_{23}=1$. The weighted degrees are $\deg_w(0)=1$, $\deg_w(1)=3$, $\deg_w(2)=3$, $\deg_w(3)=1$. We choose $\lambda=1$, which satisfies $\lambda > \Delta_w/(N+1) = 3/5$. Substituting these values into the general formulas yields the concrete QUBO parameters:
\begin{align}
    Q_{0,0} &= 1 + 1(1-4) = -2, & Q_{1,1} &= 3 + 1(1-4) = 0, \\
    Q_{2,2} &= 3 + 1(1-4) = 0, & Q_{3,3} &= 1 + 1(1-4) = -2, \\
    Q_{0,1} &= -2(1) + 2(1) = 0, & Q_{1,2} &= -2(2) + 2(1) = -2, \\
    Q_{2,3} &= -2(1) + 2(1) = 0, & Q_{0,2} &= 2(1) = 2, \\
    Q_{0,3} &= 2(1) = 2, & Q_{1,3} &= 2(1) = 2, \\
    c &= \frac{1 \cdot 4^2}{4} = 4.
\end{align}
The resulting Hamiltonian is:
\begin{align}
    H(x) = -2x_0 - 2x_3 - 2x_1 x_2 + 2x_0 x_2 + 2x_0 x_3 + 2x_1 x_3 + 4.
\end{align}
Evaluating $H(x)$ on the balanced assignment $x = (1,1,0,0)$ (partition $\{0,1\}$ vs $\{2,3\}$) gives $H(1,1,0,0) = -2(1) - 2(0) - 2(0) + 0 + 0 + 0 + 4 = 2$. This matches the cut weight of the partition, which consists solely of the edge $(1,2)$ with weight $2$. The global minimum of $H(x)$ over all balanced assignments is $2$, confirming that $\{0,1\}$ and $\{2,3\}$ form the minimum cut bisection.
\end{document}